\chapter{Firewall, Security Groups y NACL}

La solución utiliza Security Groups como firewall stateful asociado a las instancias y Network
ACL como control stateless asociado a las subredes. Esta combinación implementa defensa en
profundidad sin incorporar AWS Network Firewall, cuyo costo y complejidad no son necesarios
para el escenario académico. La aceptación de esta interpretación debe confirmarse con el
docente antes del despliegue.

\section{Política de Visitas}

El Security Group de Visitas permite DNS hacia el resolvedor de la VPC y HTTP/HTTPS hacia
Internet. La NACL evalúa primero una denegación dirigida a \code{10.10.0.0/24}; por ello los
permisos generales hacia \code{0.0.0.0/0} no abren la intranet ni otras redes internas. El
bastión conserva SSH administrativo y TI puede iniciar pruebas controladas.

\section{Política de TI}

TI puede iniciar tráfico hacia la VPC. Su Security Group solo acepta SSH desde el bastión; no
existe una regla de entrada para Ventas, Visitas, Data Center o DMZ. La NACL permite puertos
efímeros de retorno y deniega el resto del tráfico iniciado hacia el segmento.

\section{Política de DMZ}

La DMZ acepta HTTP público y HTTPS cuando la opción correspondiente está activa. SSH se
limita al Security Group del bastión. Las reglas de salida permiten DNS, actualizaciones y
respuestas a clientes, pero una denegación de menor número bloquea conexiones iniciadas hacia
la VPC antes de alcanzar los permisos generales de Internet.

\begin{warningbox}[Naturaleza stateless]
Las NACL no conservan estado. Los retornos TCP y UDP requieren reglas explícitas para puertos
efímeros \code{1024-65535}. Los Security Groups continúan siendo el control principal de
identidad y servicio en cada instancia.
\end{warningbox}

\section{Hardening incorporado en la base}

La base mantiene IMDSv2 obligatorio, EBS cifrado, SSH restringido al bastión, cómputo de Fase
2 deshabilitado por defecto y ausencia de claves o estados en Git. El hardening del sistema
operativo y de Nginx/BIND9 se integrará con los archivos de los responsables de servicio.
